
% BEGIN bib
 
\par\medskip\noindent\textbf{References}\par
  
\par\noindent  [1]  Wilhelm Ackermann.  Begründung des “tertium non datur” mittels der Hilbertschen Theorie der Widerspruchsfreiheit.  Mathematische Annalen, 93:1–36, 1924.  doi:10.1007/BF01449946.  
 
\par\noindent  [2]  Federico Aschieri and Stefano Berardi.  Interactive learning-based realizability for Heyting arithmetic with EM1.  Logical Methods in Computer Science, 6(3), 2010.  doi:10.2168/LMCS-6(3:19)2010.  
 
\par\noindent  [3]  Federico Aschieri, Adata Ciabattoni, and Francesco Genco.  On the concurrent computational content of intermediate logics.  Theoretical Computer Science, 813:375–409, 2020.  doi:10.1016/j.tcs.2020.01.022.  
 
\par\noindent  [4]  Jeremy Avigad, Edward T. Dean, and Jason Rute.  A metastable dominated convergence theorem.  Journal of Logic \& Analysis, 4(3):1–19, 2012.  doi:10.4115/jla.2012.4.3.  
 
\par\noindent  [5]  Jeremy Avigad, Philipp Gerhardy, and Henry Towsner.  Local stability of ergodic averages.  Transactions of the American Mathematical Society, 362(1):261–288, 2010.  doi:10.1090/S0002-9947-09-04814-4.  
 
\par\noindent  [6]  Stefano Berardi, Marc Bezem, and Thierry Coquand.  On the computational content of the axiom of choice.  Journal of Symbolic Logic, 63(2):600–622, 1998.  doi:10.2307/2586854.  
 
\par\noindent  [7]  Ulrich Berger and Helmut Schwichtenber.  Program extraction from classical proofs.  In Logic and Computational Complexity workshop (LCC’94), volume 960 of Lecture Notes in Computer Science, pages 77–97, 1995.  doi:10.1007/3-540-60178-3\_80.  
 
\par\noindent  [8]  Valeria de Paiva.  The Dialectica categories.  PhD thesis, University of Cambridge, 1991.  Published as Technical Report 213, Computer Laboratory, University of Cambridge.  doi:10.48456/tr-213.  
 
\par\noindent  [9]  Jerry den Hartog and Erik P. de Vink.  Verifying probabilistic programs using a Hoare like logic.  International Journal of Foundations of Computer Science, 13:315–340, 2002.  doi:10.1142/S012905410200114X.  
 
\par\noindent  [10]  Justus Diller.  Eine Variante zur Dialectica-Interpretation der Heyting-Arithmetik endlicher Typen.  Archiv für mathematische Logik und Grundlagenforschung, 16(1–2):49–66, 1974.  doi:10.1007/BF02025118.  
 
\par\noindent  [11]  Philipp Gerhardy and Ulrich Kohlenbach.  General logical metatheorems for functional analysis.  Transactions of the American Mathematical Society, 360:2615–2660, 2008.  doi:10.1090/S0002-9947-07-04429-7.  
 
\par\noindent  [12]  Kurt Gödel.  Über eine bisher noch nicht benützte Erweiterung des finiten Standpunktes.  dialectica, 12(3–4):280–287, 1958.  doi:10.1111/j.1746-8361.1958.tb01464.x.  
 
\par\noindent  [13]  Timothy Griffin.  A formulae-as-type notion of control.  In Proceedings of Principles of Programming Languages (POPL’90), pages 47–58. ACM, 1990.  doi:10.1145/96709.96714.  
 
\par\noindent  [14]  Charles Antony Richard Hoare.  An axiomatic basis for computer programming.  Communications of the ACM, 12(10):576–580, 1969.  doi:10.1145/363235.363259.  
 
\par\noindent  [15]  William A. Howard.  Hereditarily majorizable functionals of finite type.  In Anne S. Troelstra, editor, Metamathematical Investigation of Intuitionistic Arithmetic and Analysis, volume 344 of Lecture Notes in Mathematics. Springer-Verlag, 1973.  doi:10.1007/BFb0066739.  
 
\par\noindent  [16]  Marie Kerjean and Pierre-Marie Pédrot.  $\delta$ is for Dialectica.  In Proceedings of Logic in Computer Science (LICS’24), pages 48:1–13. ACM, 2024.  doi:10.1145/3661814.3662106.  
 
\par\noindent  [17]  Ulrich Kohlenbach.  Some logical metatheorems with applications in functional analysis.  Transactions of the American Mathematical Society, 357(1):89–128, 2005.  doi:10.1090/s0002-9947-04-03515-9.  
 
\par\noindent  [18]  Ulrich Kohlenbach.  Applied Proof Theory: Proof Interpretations and their Use in Mathematics.  Springer Monographs in Mathematics. Springer, 2008.  doi:10.1007/978-3-540-77533-1.  
 
\par\noindent  [19]  Ulrich Kohlenbach.  Recent progress in proof mining in nonlinear analysis.  IFCoLoG Journal of Logics and their Applications, 10(4):3361–3410, 2017.  URL: https://www2.mathematik.tu-darmstadt.de/˜kohlenbach/progress-new.pdf.  
 
\par\noindent  [20]  Ulrich Kohlenbach.  Proof-theoretic methods in nonlinear analysis.  In Proceedings of the International Congress of Mathematicians 2018, volume 2, pages 61–82. World Scientific, 2019.  doi:10.1142/9789813272880\_0045.  
 
\par\noindent  [21]  Ulrich Kohlenbach.  Local formalizations in nonlinear analysis and related areas and proof-theoretic tameness.  In Kreisel’s Interests: On the Foundations of Logic and Mathematics, volume 41 of Tributes, pages 45–61. College Publications, 2020.  URL: https://www2.mathematik.tu-darmstadt.de/˜kohlenbach/Kreisel2018.pdf.  
 
\par\noindent  [22]  Jean-Louis Krivine.  Realizability in classical logic. in: Interactive models of computation and program behaviour.  Panoramas et synthéses, 27:197–229, 2009.  URL: https://hal.science/hal-00154500/document.  
 
\par\noindent  [23]  Wilfried Meyer-Viol.  Instantial Logic.  PhD thesis, University of Amsterdam, 1985.  URL: https://eprints.illc.uva.nl/id/eprint/1984/2/DS-1995-11.text.pdf.  
 
\par\noindent  [24]  Alexandre Miquel.  Existential witness extraction in classical realizability and via a negative translation.  Logical Methods in Computer Science, 7(2:2):1–47, 2011.  doi:10.2168/LMCS-7(2:2)2011.  
 
\par\noindent  [25]  Morenikeji Neri and Nicholas Pischke.  Proof mining and probability theory.  Forum of Mathematics, Sigma, 13:e187, 2025.  doi:10.1017/fms.2025.10138.  
 
\par\noindent  [26]  Morenikeji Neri, Nicholas Pischke, and Thomas Powell.  Generalized learnability of stochastic principles.  In Proceedings of Computability in Europe (CiE’25), volume 15764 of LNCS, pages 333–348, 2025.  doi:10.1007/978-3-031-95908-0\_24.  
 
\par\noindent  [27]  Morenikeji Neri and Thomas Powell.  On quantitative convergence for stochastic processes: Crossings, fluctuations and martingales.  Transactions of the American Mathematical Society, Series B, 12:974–1019, 2025.  doi:10.1090/btran/229.  
 
\par\noindent  [28]  Morenikeji Neri and Thomas Powell.  A quantitative Robbins-Siegmund theorem.  Annals of Applied Probability, 2025.  To appear.  doi:10.48550/arXiv.2410.15986.  
 
\par\noindent  [29]  Peter O’Hearn and David Pym.  The logic of bunched implications.  Bulletin of Symbolic Logic, 5(2):215–244, 1999.  doi:10.2307/421090.  
 
\par\noindent  [30]  Paulo Oliva.  Unifying functional interpretations.  Notre Dame Journal of Formal Logic, 47(2):263–290, 2006.  doi:10.1305/ndjfl/1153858651.  
 
\par\noindent  [31]  Paulo Oliva.  Functional interpretations of linear and intuitionistic logic.  Information and Computation, 208(5):565–577, 2010.  doi:10.1016/j.ic.2008.11.008.  
 
\par\noindent  [32]  Michel Parigot.  $\lambda\mu$-calculus: An algorithmic interpretation of classical natural deduction.  In Proceedings of Logic Programming and Automated Reasoning (LPAR’92), pages 190–201. Springer Berlin Heidelberg, 1992.  doi:10.1007/BFb0013061.  
 
\par\noindent  [33]  Pierre-Marie Pédrot.  A functional functional interpretation.  In Joint proceedings of Computer Science Logic and Logic in Computer Science (CSL–LICS’14), pages 77:1–10. ACM, 2014.  doi:10.1145/2603088.2603094.  
 
\par\noindent  [34]  Thomas Powell.  Gödel’s functional interpretation and the concept of learning.  In Proceedings of Logic in Computer Science (LICS ’16), pages 136–145. ACM, 2016.  doi:10.1145/2933575.2933605.  
 
\par\noindent  [35]  Thomas Powell.  Computational interpretations of classical reasoning: From the epsilon calculus to stateful programs.  In Mathesis Universalis, Computability and Proof, volume 412 of Synthese Library, pages 255–290. Springer, 2019.  doi:10.1007/978-3-030-20447-1\_14.  
 
\par\noindent  [36]  Thomas Powell.  Proofs as stateful programs: A first-order logic with abstract Hoare triples.  Logical Methods in Computer Science, 20(1):7:1–7:32, 2024.  doi:10.46298/LMCS-20(1:7)2024.  
 
\par\noindent  [37]  John C. Reynolds.  Separation logic: a logic for shared mutable data structures.  In Proceedings of Logic in Computer Science (LICS’02), pages 55–74. IEEE, 2002.  doi:10.1109/LICS.2002.1029817.  
 
\par\noindent  [38]  Helmut Schwichtenberg.  Dialectica interpretation of well-founded induction.  Mathematical Logic Quarterly, 54(3):229–239, 2008.  doi:10.1002/malq.200710045.  
 
\par\noindent  [39]  Davide Trotta, Jonathan Weinberger, and Valeria de Paiva.  Skolem, Gödel, and Hilbert fibrations.  available at https://arxiv.org/abs/2407.15765, 2024.  
  
% END bib
